% from ega2-1-fr.tex
\oldpage[II]{189}
et on désigne par $V$ le schéma déduit de $C'$ par éclatement de $b$~; $X$ est
le sous-préschéma fermé de $V$, image réciproque de $a$ par le morphisme
structural $q:V\to C$. On montre que $X$ est réunion de deux composantes
irréductibles $X_1$, $X_2$, où $X_1$ est l'image réciproque de $b$ dans $V$.
Il est immédiat que l'Idéal $\mathcal{J}$ de $\mathcal{O}_V$ qui définit $X$
est encore inversible, mais on prouve que
$j^*(\mathcal{J})=\mathcal{L}$ ($j$ injection canonique $X\to V$) n'est pas
ample, par la considération du «~degré~» de l'image réciproque de
$\mathcal{L}$ sur $X_1$, qui devrait être $>0$ si $\mathcal{L}$ était ample,
et qu'on montre (par un calcul élémentaire d'intersections) être égal à 0.
\end{remark}

\subsection{Unicité des contractions}
\label{subsection:II.8.11-fr}

\begin{lemma}[8.11.1]
\label{II.8.11.1-fr}
Soient $U$, $V$ deux préschémas,
$h=(h_0,\lambda):U\to V$ un morphisme surjectif. On suppose que~:
\begin{enumerate}
  \item[$1^\circ$]
  $\lambda:\mathcal{O}_V\to h_*(\mathcal{O}_U)
  =(h_0)_*(\mathcal{O}_U)$ est un isomorphisme.
  \item[$2^\circ$] L'espace sous-jacent à $V$ s'identifie à l'espace quotient
  de l'espace sous-jacent à $U$ par la relation
  $R:h_0(x)=h_0(y)$ (condition toujours vérifiée lorsque le morphisme $h$ est
  ouvert, ou fermé, ou a fortiori lorsque $h$ est propre).
\end{enumerate}
Alors, pour tout préschéma $W$, l'application
\begin{equation}
\label{II.8.11.1.1-fr}
  \operatorname{Hom}(V,W)\longrightarrow\operatorname{Hom}(U,W)
\tag{8.11.1.1}
\end{equation}
qui, à tout morphisme $v=(v_0,\nu)$ de $V$ dans $W$, fait correspondre le
morphisme $u=v\circ h=(u_0,\mu)$, est une bijection de
$\operatorname{Hom}(V,W)$ sur l'ensemble des $u$ tels que $u_0$ soit
constante sur toute fibre $h_0^{-1}(x)$.
\end{lemma}

\begin{proof}
Il est clair que si $u=v\circ h$, donc $u_0=v_0\circ h_0$, $u_0$ est
constante sur tout ensemble $h_0^{-1}(x)$. Inversement, si $u$ a cette
propriété, montrons qu'il existe un $v\in\operatorname{Hom}(V,W)$ et un seul
tel que $u=v\circ h$. L'existence et l'unicité de l'application continue
$v_0:V\to W$ telle que $u_0=v_0\circ h_0$ résultent de l'hypothèse, puisque
$h_0$ s'identifie à l'application canonique de $U$ sur $U/R$. On peut d'autre
part, en remplaçant au besoin $V$ par un préschéma isomorphe, supposer que
$\lambda$ est l'identité~; par hypothèse, $\mu$ est alors un homomorphisme
\[
  \mu:\mathcal{O}_W\longrightarrow(u_0)_*(\mathcal{O}_U)
  =(v_0)_*((h_0)_*(\mathcal{O}_U))
\]
tel que l'homomorphisme correspondant
$\mu^\sharp:u_0^*(\mathcal{O}_W)\to\mathcal{O}_U$ soit local sur chaque
fibre. Comme
$(v_0)_*((h_0)_*(\mathcal{O}_U))=(v_0)_*(\mathcal{O}_V)$, on doit
nécessairement avoir $\nu=\mu$, et tout revient à voir que l'homomorphisme
correspondant $\nu^\sharp:v_0^*(\mathcal{O}_W)\to\mathcal{O}_V$ est local sur
chaque fibre. Or, tout $y\in V$ est de la forme $h_0(x)$ pour un $x\in U$~;
posons $z=v_0(y)=u_0(x)$. Alors
(\hyperref[0.3.5.5-fr]{0, 3.5.5}) l'homomorphisme $\mu_x^\sharp$ se factorise
en
\[
  \mu_x^\sharp:\mathcal{O}_z
  \xrightarrow{\nu_y^\sharp}\mathcal{O}_y
  \xrightarrow{\lambda_x^\sharp}\mathcal{O}_x.
\]
Par hypothèse $\lambda_x^\sharp$ et $\mu_x^\sharp$ sont des homomorphismes
locaux~; donc $\lambda_x^\sharp$ transforme tout élément inversible de
$\mathcal{O}_y$ en un élément inversible de $\mathcal{O}_x$~; si
$\nu_y^\sharp$ transformait un élément non inversible de $\mathcal{O}_z$ en
un élément inversible de $\mathcal{O}_y$, $\mu_x^\sharp$ transformerait cet
élément de $\mathcal{O}_z$ en un élément inversible de $\mathcal{O}_x$,
contrairement à l'hypothèse, d'où le lemme.
\end{proof}

\begin{corollary}[8.11.2]
\label{II.8.11.2-fr}
Soient $U$ un préschéma intègre, $V$ un préschéma normal~; alors tout morphisme
$h:U\to V$ qui est universellement fermé, birationnel et radiciel est un
isomorphisme.
\end{corollary}

\begin{proof}
Si $h=(h_0,\lambda)$, il résulte des hypothèses que $h_0$ est injectif et
fermé, et que $h_0(U)$

